\scanpage{28}

En vertu \unsure{du critère|des axiomes|on sait} qui \unsure{suit}, si $A$ est la sous-catégorie pleine de $C$ engendrée par les $\Phi(i,n)$ et les $\Psi^i$, alors $C$ est connue à équiv près (\unsure{déf à isom.|déf. isom.} unique près) si on connaît $A$ avec sa structure $\frac12$-additive. Or c'est \inserted{elle} qui vient d'être déterminée (on a omis de décrire la composition des morphismes, qui est évidente \ldots).

On trouve donc que $C\simeq$ groupes $\frac12$ simpliciaux abéliens str. $\partial$-parfaits \ldots.

La chose restant à faire est d'en déterminer la $\otimes$-structure et les opérations $\Gamma^{\unsure{i|\prime}}$.

% archive sheet number: 27
